\section{A Worked Example: Slow Poisoning in a Hospital Ward}\label{sec:example}

This illustrative scenario (not an experiment) makes the problem concrete. A ward network carries an
infusion pump, a patient monitor and a smart bed from three different makers. Each runs an anomaly
detector that keeps learning from its own traffic. An attacker has taken control of the pump through a
firmware vulnerability. The goal is to send patient data to the attacker's server, or receive commands
from it, without an alarm.

\textbf{The slow poisoning.} Table~\ref{tab:frog} shows the campaign. Each step is small enough not to
alarm, so each step is learned, and the detector's idea of normal creeps toward the exfiltration traffic.
When the large transfer finally happens it is inside the learned normal region. The pump's detector is
\emph{compromised}: it now misses an attack it would have caught before.

\begin{table}[t]
\caption{A slow poisoning campaign against a hospital infusion pump (illustrative)}\label{tab:frog}
\centering\footnotesize
\begin{tabular}{@{}p{1.2cm}p{3.6cm}p{3.0cm}@{}}
\toprule
Time & Pump traffic & Lone detector \\
\midrule
Week 1 & A few bytes per hour to a new server & Below threshold, learned \\
Week 2 & Slightly larger, slightly more often & Below the moved threshold, learned \\
Weeks 3--6 & Steadily larger transfers & Each step learned \\
Attack & Bulk transfer of patient records & Inside learned normal: \emph{no alarm} \\
\bottomrule
\end{tabular}
\end{table}

\textbf{Why the lone detector cannot tell.} A real firmware update could produce the same slow change.
From the pump's traffic alone, the honest story and the attack story are indistinguishable.

\textbf{What FOA does.} Each cycle the pump asks its neighbours, in effect, ``do you see this change
too?'' The monitor and the bed do not: their traffic is unchanged, so the consensus shows no shared
change (low significance $w$) and the pump's own core traffic does not move along any shared direction
(low agreement $\lambda$). The change is therefore device-only and passes through the residual budget
$\norm{R(t)}\le(B_0+Bt)\rho$, which starts at a tenth of an alarm radius and grows by 0.005 per cycle.
By Corollary~\ref{cor:ttc}, an exfiltration target two alarm radii away cannot be hidden before about
133 cycles, about five and a half days at one learning cycle per hour, and targets further away take
proportionally longer. Throughout that time the pump keeps alarming on the attacker's larger transfers,
which gives staff time to notice.

\textbf{The honest contrast.} When the hospital moves the ward to a new network gateway, all three
devices' traffic changes together. The neighbours back the change, $w$ and $\lambda$ rise, and FOA
learns it through the larger common budget, without the flood of false alarms a frozen detector would
raise.

\textbf{The same effect in our data.} Figure~\ref{fig:frog} shows the classic boundary poisoner in
simulation on real N-BaIoT traffic during strong shared drift. A lone device and the DriftXpert-style
gate are walked down to about 16\% detection of the target attack; FOA keeps detecting it, with victim
false alarms near 3\%.
